\pdfoutput=1
\documentclass[11pt]{article}
\usepackage[margin=1in]{geometry}
\usepackage{amsmath,amssymb,amsthm,booktabs}
\usepackage{hyperref}
\hypersetup{hidelinks,
 pdftitle={Two constants in the numerical estimates of inter-universal Teichmuller theory IV},
 pdfauthor={Toshihisa Urahigashi}}
\newtheorem{proposition}{Proposition}
\theoremstyle{remark}
\newtheorem{remark}{Remark}
\newcommand{\Fmod}{F_{\mathrm{mod}}}
\newcommand{\ord}{\operatorname{ord}}
\newcommand{\qpilot}{\underline{\underline{q}}}
\newcommand{\thetapilot}{\underline{\underline{\Theta}}}
\newcommand{\Smark}{\textup{[S]}}
\newcommand{\Pmark}{\textup{[P]}}
\title{Two constants in the numerical estimates of\\
 inter-universal Teichm\"uller theory IV\\[4pt]
 \large $12/(\ell(\ell+1))$ and $12/\ell^2$: a fifth report}
\author{Toshihisa Urahigashi\\
 \small\texttt{123026.urahigashi.toshihisa@gmail.com}}
\date{Version 1.0.0 --- 8 September 2026\\[3pt]
 \small ORCID \href{https://orcid.org/0009-0004-0460-6242}{0009-0004-0460-6242}}

\begin{document}
\maketitle
\begin{abstract}
The constants $\chi=12/(\ell(\ell+1))$ and $\tau=12/\ell^2$
occur at different steps of the numerical estimates surrounding IUT IV.
The first is the exact ratio of the stage-one coefficient $1/(2\ell)$ to the
mean of $j^{2}/(2\ell)$ over $j=1,\dots,(\ell-1)/2$, and appears in
Scholze--Stix's approximate calculation.  The second is used in the upper-bound rearrangement in the
proof of IUT IV, Theorem~1.10.  We locate both in the original texts and
compute their difference with exact rational arithmetic.  With all other
inputs fixed, the latter rearrangement adds $H/(2\ell^2)$ to the unnormalized
expression, or $1/\ell$ after division by $H/(2\ell)$; its reciprocal factor
is strictly larger.  Tables at $\ell=7,157$ also give the contributions from
the designated split primes of two explicit curves.  These curves serve only
as numerical illustrations; none of the algebraic propositions depends on
their initial-data admissibility.  What is compared here are two scalar expressions under a common set of inputs;
this is not a comparison of the two sources' full estimates, whose positive
terms, divisor supports and approximation conventions differ and are kept
separate throughout.  These calculations do not determine a possible-image hull
or adjudicate IUT or the Scholze--Stix objection.
\end{abstract}

\section{The objects and the comparison being made}\label{s:scope}

The coefficient identity in \cite[\S7]{Fourth} is the starting point.
Here we identify a larger sufficient constant used in an upper-bound
estimate, compare it with the sharp threshold for the scalar rearrangement
defined below, and quantify
the difference both before and after the source's root normalization.
The finite-sum identity is elementary and is not claimed as a new discovery.
Nor is priority claimed for comparing the two constants: the purpose is
to record the exact source steps, common-input conditions and numerical
differences in one place.

We reserve \Smark{} for the explicitly quoted text or mathematical display
of a source; \Pmark{} marks the deductions proved here, including the
comparison of the displayed expressions.  A source assertion and a proof
of that assertion are different things.  The hypotheses of the algebraic
comparison below are written out; their satisfaction by an actual
possible-image hull is not supplied by the arithmetic checks.
All references to pages of IUT I and IV are to the consulted author
versions, May 2020 and April 2020 respectively.  Publication metadata are
listed separately.  For Scholze--Stix we use the manuscript of July 16, 2018.

Write $H$ for a normalized degree of an \emph{unrooted} Tate-parameter
divisor and set
\begin{equation}\label{e:constants}
 m=\frac{\ell-1}{2},\qquad r=\frac1{2\ell},\qquad
 Q=rH,\qquad b=\frac{\ell+1}{24},\qquad c=\frac{\ell+1}{4}.
\end{equation}
Until a source application is specified, these are real or rational
quantities, with $\ell$ an odd integer at least~5 and $H>0$.
In IUT IV the relevant $H$ is $\log(\mathfrak q)$, the normalized degree
of the divisor supported on the \emph{selected} bad places
\cite[Thm.~1.10, pp.~22--23]{IUT4}.  The prose normalization there gives the following equality in our notation:
\begin{equation}\label{e:root}
 \Pmark\qquad |\log(\qpilot)|=\frac1{2\ell}\log(\mathfrak q)
 \qquad\text{\cite[Thm.~1.10, p.~23]{IUT4}}.
\end{equation}
Here the double-underlined $\qpilot$ is the source's pilot notation;
it is not an additional Tate parameter to which a second root should be applied.

Scholze--Stix define $q_E$ by summing over \emph{all} finite places of
multiplicative reduction, and use $\deg(q_E)/(2\ell)$ for the rooted
quantity \cite[\S1.2, p.~2; \S1.3, pp.~3--4]{SS}.
Consequently, using the same elliptic curve is not alone enough to identify
their $\deg(q_E)$ with IUT IV's selected $H$.
Equality requires the same support, or explicit accounting for the omitted
contribution.  Under the extra good-reduction hypothesis of Theorem~1.10,
the selected and actual bad places agree outside $2\ell$; the contribution
over $2\ell$ must still be checked.  Likewise the positive terms appearing
in the two texts are not identified merely by giving them the same symbol.

\section{What is printed in Scholze--Stix}\label{s:ss}

Their setup in \S1.2 fixes a degree bound for $k$, restricts the local data
at $2$ and infinity, and uses split semistable reduction.  The prime
$\ell=\ell(P)$ is also constrained by the height window (1.3), the
mod-$\ell$ image, and the local Tate extension condition; an exceptional
list is set aside \cite[\S1.2, pp.~2--3]{SS}.
The substitutions in \S\ref{s:tables} evaluate coefficients, and do not
assert these additional conditions for either displayed level.

\Smark{} Their (1.4), p.~3 \cite[p.~3, (1.4); p.~4, (1.5)]{SS}, is
\begin{equation}\label{e:ss14}
 -|\log(\qpilot)|\le -|\log(\thetapilot)|.
\end{equation}
On p.~4 the introduction to (1.5) reads
``The right side of (1.4) is in essential approximation''.
The immediately following parenthesis defines $\mathfrak d(P)_v$ as the
$v$-part of the logarithmic root discriminant.  Retaining the source's
$\ell^*$, the displayed average is
\begin{equation}\label{e:ss15}
 \Smark\qquad
 \begin{aligned}
 -|\log(\thetapilot)|&\approx
 \frac1{\ell^*}\sum_v\sum_{j=1}^{\ell^*}
 \left(j\cdot\mathfrak d(P)_v-
 j^2\frac1{2\ell[k:\mathbb Q]}v(q_v)\log N(v)\right)
 \\&=\frac{\ell^*+1}{2}\left(\mathfrak d(P)-\frac16\deg(q_E)\right).
 \end{aligned}
\end{equation}
The intermediate sum-of-squares evaluation in their display is omitted;
no approximation sign is replaced by an equality.  In our deductions
$m=\ell^*=(\ell-1)/2$.  The following \emph{unnumbered} display is
\begin{equation}\label{e:ssfinal}
 \Smark\qquad
 \frac16\deg(q_E)\lessapprox
 \left(1-\frac{12}{\ell(\ell+1)}\right)^{-1}\cdot\mathfrak d(P).
\end{equation}
The next sentence says the error is handled
``with a linear term in $\ell$ on the right side''.
It then states that this achieves Claim~6; the following sentence begins
the authors' objection to the reasoning for (1.5), proposing (1.6) with
the $j^2$ factor missing.  We are quoting the approximate calculation
they describe, not attributing their endorsement to it.
Thus \eqref{e:ssfinal} must not be quoted as an
error-free inequality with $\le$, nor called their equation~(1.5).
The exact bound proved below is instead conditional on an explicitly
specified additive error.

\begin{proposition}[the two finite moments; \Pmark]\label{p:moments}
For odd $\ell\ge5$,
\[
 \frac1m\sum_{j=1}^m j=c,\qquad
 \frac1m\sum_{j=1}^m\frac{j^2}{2\ell}=b=\frac c6,
 \qquad \frac br=\frac{\ell(\ell+1)}{12}.
\]
In particular the unique number $t$ for which $bt=r$ is
\begin{equation}\label{e:chi}
 \chi=\frac rb=\frac{12}{\ell(\ell+1)},\qquad
 A=\frac1\chi=\frac{\ell(\ell+1)}{12}.
\end{equation}
\end{proposition}
\begin{proof}
The sums are $m(m+1)/2$ and $m(m+1)(2m+1)/6$.
Substituting $m=(\ell-1)/2$ gives $c=(\ell+1)/4$ and
$b=(m+1)(2m+1)/(12\ell)=(\ell+1)/24$.
Dividing $r$ by this positive $b$ gives \eqref{e:chi}.
\end{proof}

This explains both $(\ell+1)/4$ and $1/6$ in the algebra of
\eqref{e:ss15}: $c$ is the first moment and $1/6=b/c$.
SS also discuss the geometric origin of $1/6$ in their height comparison
(1.2), p.~2 \cite[p.~2, (1.2)]{SS}: the respective degrees there are $1/12$ and $1/2$.
That discussion is subject to their bounded local data; it does not assert
an exact equality between every height and $\deg(q_E)/6$.
Their displayed sums and rearrangement account for the denominator
$\ell(\ell+1)$.  Our comparison with the $\ell^2$ denominator uses the
specific IUT IV estimate below, without attributing an unstated motive
or claiming a first occurrence of this comparison.

\section{The sufficient constant used in IUT IV}\label{s:iv}

\Pmark{} Reading the negative terms in Theorem~1.10, proof, Step~(v), p.~29
\cite[Thm.~1.10, proof, Step~(v), p.~29]{IUT4}, gives the following equality of
negative terms in the first two lines of its procession-normalized upper
bound (with $\ell^*=(\ell-1)/2$ as in the source):
\[
 -\frac{(2\ell^*+1)(\ell^*+1)}{12\ell}\log(\mathfrak q_{v_{\mathbb Q}})
 =-\frac{\ell+1}{24}\log(\mathfrak q_{v_{\mathbb Q}}).
\]
The same display has different term
$\frac{\ell+5}{4}\log(\mathfrak d^K_{v_{\mathbb Q}})$.
This is the average of $j+1$, whereas \eqref{e:ss15} averages $j$.
\textup{[S]} \emph{The same rearrangement, with the same slack, is in print
elsewhere.}  \cite[p.~70]{JoshiIV} writes: ``Hence one may replace the term
$-\frac16\cdot\log(q)$ in (6.11.3) by $-\frac16\cdot\bigl(1-\frac{12}{\ell^2}\bigr)\cdot
\log(q)$ and still have an excess contribution greater than
$\frac1{2\ell}\log(q)$.''  No priority is claimed here for the rearrangement or
for its slack; what this report adds is the exact comparison with $\chi$ of
\S\ref{s:ss}.  \emph{This report takes no position on the further claims of}
\cite{JoshiIV}.

The next upper bound factors out $(\ell+1)/4$ and replaces the remaining
positive coefficients by larger ones.  For example,
\[
 \Pmark\qquad \frac{\ell+5}{4}
 =\frac{\ell+1}{4}\left(1+\frac4{\ell+1}\right)
 \le\frac{\ell+1}{4}\left(1+\frac4\ell\right).
\]
Thus the common negative coefficient does not make the full positive
parts of the two sources equal.

The occurrence of $12/\ell^2$ can be located independently of this
interpretation.  \Smark{} The statement of Theorem~1.10, p.~23
\cite[Thm.~1.10, p.~23; proof, Step~(viii), p.~30]{IUT4}, contains
$-\frac16(1-12/\ell^2)\log(\mathfrak q)$ in its displayed choice of $C_\Theta$.
The proof, Step~(viii), p.~30, ends the explanation of that upper-bound
display with the complete estimate clause:
\begin{quote}
\Smark{} ``--- where we apply the estimate
$\frac{\ell+1}{4}\cdot\frac16\cdot\frac{12}{\ell^2}\ge\frac1{2\ell}$
[cf.\ the fact that $\ell\ge1$].''
\end{quote}
Let
\[
 \tau=\frac{12}{\ell^2},\qquad
 u=b\tau=\frac{\ell+1}{2\ell^2},\qquad
 \epsilon=u-r=\frac1{2\ell^2}.
\]
The following identity pinpoints the change in the upper bound.

\begin{proposition}[the rearrangement and its slack; \Pmark]\label{p:slack}
For $H>0$ and any real $D$, define the numerical expression
\[
 K_t(H,D)=c\left(D-\frac{1-t}{6}H\right)-rH.
\]
Then
\begin{align*}
 K_\chi(H,D)&=cD-bH,\\
 K_\tau(H,D)&=cD-bH+\frac{H}{2\ell^2},\\
 \frac{K_\tau(H,D)-K_\chi(H,D)}{Q}&=\frac1\ell.
\end{align*}
Among $t\in\mathbb R$, the inequality
$K_t(H,D)\ge cD-bH$ holds if and only if $t\ge\chi$.
\end{proposition}
\begin{proof}
Subtracting $cD-bH$ leaves $(bt-r)H$.
At $t=\chi$ it is zero.  At $t=\tau$ it is
$((\ell+1)/(2\ell^2)-1/(2\ell))H=H/(2\ell^2)$.
Division by $Q=H/(2\ell)$ gives $1/\ell$.
Since $b,H>0$, its sign is the sign of $t-\chi$.
\end{proof}

In Step~(viii) the terms represented by $D$ include different, conductor
and other positive estimates; a further prime-counting estimate follows
on pp.~30--31.  Proposition~\ref{p:slack} compares the expressions with
\emph{these other terms held fixed}.  In particular, $1/\ell$ is the
change of this normalized numerical expression, not a measured change
in the actual $C_\Theta$ or in a possible-image hull.

\section{Comparison of the two reciprocal scalar factors}\label{s:order}

The comparison below is between two \emph{scalar expressions} evaluated on a
common set of inputs, which Proposition~\ref{p:bound} states as a hypothesis.
It is not a comparison of the full estimates of \cite{SS} and \cite{IUT4}: no
map identifying their supports, positive terms, error terms or approximation
conventions is given here, and none is claimed.

\begin{proposition}[an explicit error and common inputs; \Pmark]\label{p:bound}
Let $H>0$, let $D,R$ be real, and put $U=D+R/c\ge0$.
Suppose a real number $h$ satisfies the two bounds
\[
 -rH\le h\le c(D-H/6)+R.
\]
Then
\begin{equation}\label{e:strictbound}
 \frac H6\le F_\chi U\le F_\tau U,
 \qquad F_\chi=\frac1{1-\chi},\quad F_\tau=\frac1{1-\tau}.
\end{equation}
The second inequality is strict when $U>0$.  The exact differences are
\begin{align}
 \tau-\chi&=\frac{12}{\ell^2(\ell+1)},&
 \frac\tau\chi-1&=\frac1\ell,\label{e:difftau}\\
 F_\tau-F_\chi&=
 \frac{12\ell}{(\ell^2-12)(\ell(\ell+1)-12)},&
 \frac{F_\tau}{F_\chi}-1&=
 \frac{12}{(\ell+1)(\ell^2-12)}.\label{e:difffactor}
\end{align}
\end{proposition}
\begin{proof}
The two assumed bounds give $(b-r)H\le cD+R$.
As $b=c/6$ and $r=b\chi$, division by $c$ gives
$(1-\chi)H/6\le U$.  For odd $\ell\ge5$,
$0<\chi<\tau<1$, so $0<F_\chi<F_\tau$.
The displayed differences follow by subtraction and division of fractions.
\end{proof}

The hypothesis involving $h$ is an explicit conditional comparison;
the approximate symbol in SS is not silently replaced by it.
For nonnegative common $U$, the $\chi$ bound is tighter.
This orders the two \emph{scalar forms}, not the complete inequalities of
the two sources, whose positive terms and errors have not been equated.
Although $\tau$ and $\chi$ have the same leading order $12/\ell^2$,
their difference has order $\ell^{-3}$.  The same is true of the absolute
and relative differences of the reciprocal factors in
\eqref{e:difffactor}.  The normalized expression difference in
Proposition~\ref{p:slack} instead has order $\ell^{-1}$.
These are different comparisons, with different denominators.

\subsection*{The permitted levels}

The algebra above requires positive reciprocal denominators.
$1-\chi>0$ is equivalent to $\ell(\ell+1)>12$, and $1-\tau>0$
to $\ell^2>12$.  At the odd prime $\ell=3$ the first denominator
vanishes and the second is negative.  This is a boundary of the algebra,
not a source counterexample.

\Smark{} IUT I, Definition~3.1(c), p.~62 \cite[Def.~3.1(c), p.~62]{IUT1}, begins:
``$\ell$ is a prime number $\ge5$ such that the image of the outer
homomorphism $G_F\to GL_2(\mathbb F_\ell)$ determined by the
$\ell$-torsion points of $E_F$ contains the subgroup
$SL_2(\mathbb F_\ell)\subseteq GL_2(\mathbb F_\ell)$;''
The same item also requires $\ell$ to be prime to the selected residue
characteristics and Tate orders.  IUT IV, Theorem~1.10,
p.~22, additionally assumes full $15$-torsion over $F$ and states
$\ell\ne5$.  \Pmark{} The lower bound and primality therefore give $\ell\ge7$.
Indeed the proof's terminal estimates on p.~31 use that range and
$d_{\rm mod}=[\Fmod:\mathbb Q]\ge1$
\cite[Thm.~1.10, proof, Step~(viii), p.~31]{IUT4}:
\[
 \Smark\qquad
 (1-12/\ell^2)^{-1}\le2,\qquad
 (1-12/\ell^2)^{-1}(1+12d_{\rm mod}/\ell)
 \le1+20d_{\rm mod}/\ell.
\]
These are the source's two estimates with its stated restrictions
$\ell\ge7$ and $d_{\rm mod}\ge1$.
At $\ell=5,d_{\rm mod}=1$ the second estimate fails, since its left side
is $85/13>5$.  Level~5 can illustrate the elementary formulas but not
this theorem's terminal estimate.  The range $\ell\ge7$ by itself
does not establish the remaining initial-data, reduction or lattice
hypotheses of that theorem, including its precise torsion-generated
choice of $F$ on p.~22.

\section{Exact values at two levels}\label{s:tables}

All entries below are reduced rational numbers.  The script compares
the printed rows with independent finite-sum evaluations and the identities
proved above.  It does not convert an arithmetic parameter into admissible
initial data.

\begin{center}\small\renewcommand{\arraystretch}{1.8}\everymath{\displaystyle}
\begin{tabular}{rrrrrrr}
\toprule
$\ell$ & $\chi$ & $\tau$ & $A$ & $r$ & $u$ & $\epsilon$\\
\midrule
% BEGIN EXACT TABLE coefficients
$7$ & $\frac{3}{14}$ & $\frac{12}{49}$ & $\frac{14}{3}$ & $\frac{1}{14}$ & $\frac{4}{49}$ & $\frac{1}{98}$ \\
$157$ & $\frac{6}{12403}$ & $\frac{12}{24649}$ & $\frac{12403}{6}$ & $\frac{1}{314}$ & $\frac{79}{24649}$ & $\frac{1}{49298}$ \\
% END EXACT TABLE coefficients
\bottomrule
\end{tabular}
\end{center}

\begin{center}\small\renewcommand{\arraystretch}{1.8}\everymath{\displaystyle}
\begin{tabular}{rrrrr}
\toprule
$\ell$ & $F_\chi$ & $F_\tau$ & $F_\tau-F_\chi$ & $F_\tau/F_\chi-1$\\
\midrule
% BEGIN EXACT TABLE factors
$7$ & $\frac{14}{11}$ & $\frac{49}{37}$ & $\frac{21}{407}$ & $\frac{3}{74}$ \\
$157$ & $\frac{12403}{12397}$ & $\frac{24649}{24637}$ & $\frac{942}{305424889}$ & $\frac{6}{1946323}$ \\
% END EXACT TABLE factors
\bottomrule
\end{tabular}
\end{center}

To connect these coefficients with explicit curves, use the first report's
initial datum \cite[Thm.~1 and the arithmetic conditions]{First}:
\[
 L=\mathbb Q(\sqrt5),\quad a=(16+3\sqrt5)^{29},\quad
 E:y^2=x(x-1)(x-a),\quad\ell=7,\quad p=211.
\]
Its fields are $F=L(i,E[15])$ and $K=F(E[7])$.
For comparison, the second report's section entitled
\emph{A curve carrying fourteen candidate levels} \cite[\S13]{Second} uses
\[
 L=\mathbb Q(i),\quad a=(41+26i)^{472},\quad
 E:y^2=x(x-1)(x-a),\quad p=2357,
\]
and includes $\ell=157$ among the candidates, with
$F=L(E[15])$, $K=F(E[157])$.
The latter section explicitly distinguishes candidate levels from a
verification of all initial-data conditions at each level.
We retain that distinction: the second row below is a coefficient
evaluation at that explicit curve and level, not a new admissibility theorem.

In both designated quadratic fibres, one place is multiplicative and one
good.  At $211$, the multiplicative reduction becomes split upon the
unramified extension adjoining $i$, which leaves the Tate order unchanged;
at $2357$, it is already split over the indicated quadratic field.
If $N=\ord_{\pi}(a)$, the unrooted Tate order at the bad place is
$2N$, with $\ord_p(p)=1$.  Each base-field place has weight $1/2$, so
\begin{equation}\label{e:fibre}
 H_p=N\log p,\qquad Q_p=Nr\log p,\qquad
 (K_\tau-K_\chi)_p=N\epsilon\log p.
\end{equation}
These are contributions from that rational prime only.
No equality of the total selected degree with $H_p$ is asserted.
In particular, an extension's ramification index is not a second factor
to insert into \eqref{e:fibre}: the normalized degree is invariant under
base extension \cite[Def.~1.9(i), pp.~21--22]{IUT4}.

\begin{center}\small\renewcommand{\arraystretch}{1.8}\everymath{\displaystyle}
\begin{tabular}{rrrrrrr}
\toprule
$\ell$ & $p$ & $N$ & $Q_p/\log p$ & $bH_p/\log p$ & $uH_p/\log p$ & $\epsilon H_p/\log p$\\
\midrule
% BEGIN EXACT TABLE data
$7$ & $211$ & $29$ & $\frac{29}{14}$ & $\frac{29}{3}$ & $\frac{116}{49}$ & $\frac{29}{98}$ \\
$157$ & $2357$ & $472$ & $\frac{236}{157}$ & $\frac{9322}{3}$ & $\frac{37288}{24649}$ & $\frac{236}{24649}$ \\
% END EXACT TABLE data
\bottomrule
\end{tabular}
\end{center}

For example, the absolute difference of the reciprocal factors at
$\ell=157$ is $942/305424889$, while the designated-prime contribution to
the unnormalized expression difference is $(236/24649)\log2357$.
After division by that fibre's $Q_p$, the latter is exactly $1/157$.
A small difference of reciprocal factors and this normalized difference
are not contradictory measurements; they answer different questions.
None is an actual volume of a union of possible images.

\section*{Reproducibility and limits}

Run \texttt{python3 scripts/constant\_comparison.py} from the report bundle.
It uses the standard-library \texttt{Fraction} type, evaluates both finite
moments, checks the algebra and signs, and verifies all six printed table
rows against the TeX source.  It also requires the supplied frozen output
file and checks exact agreement with the newly computed output.
A successful run ends with
\texttt{RESULT: PASS} and \texttt{EXIT:0}; failure returns a nonzero exit
status.  Running with Python optimization enabled is rejected.
Finite numerical checks supplement the proofs; they do not prove the
identification of source objects or compute omitted source error terms.

The exact comparison orders two scalar expressions under common inputs.
It does not remove the approximation in SS, identify all its positive
terms with IUT IV's, or apply a global lower bound to a single prime.
In comments on the August 2018 SS manuscript, Mochizuki distinguishes
the linear relationships between
copies of $\mathbb R$ associated with prime-strips via the $\Theta$-link
from relationships between the log-volumes of regions before and after
indeterminacies, which he describes as dependent on the region's geometry.
For the latter relationship he writes:
\Smark{} ``In particular, this relationship is highly non-linear.''
\cite[(C14), p.~4]{Cmt}.
His report makes the same distinction and illustrates it by a region and
its symmetrization \cite[\S12, (MlLV), (LVEx), pp.~24--26]{Rpt}.
These are attributed claims about the geometric comparison, not consequences
of the scalar propositions above.  Those propositions do not identify that
comparison with multiplication by either reciprocal factor.

No conclusion is drawn about Corollary~3.12, the correctness of IUT, the
validity of the SS objection, or $abc$.
Large language models assisted drafting, source comparison and computations;
the checks reported here are internal checks.

\begin{thebibliography}{Fourth}
\bibitem[SS]{SS} P.~Scholze and J.~Stix,
\emph{Why $abc$ is still a conjecture}, manuscript, July 16, 2018, 10~pp.
\url{https://www.math.uni-bonn.de/people/scholze/WhyABCisStillaConjecture.pdf}.
\bibitem[IUT1]{IUT1} S.~Mochizuki,
\emph{Inter-universal Teichm\"uller theory I}, consulted author version,
May 2020. Published in Publ. Res. Inst. Math. Sci. \textbf{57} (2021), 3--207,
\href{https://doi.org/10.4171/PRIMS/57-1-1}{doi:10.4171/PRIMS/57-1-1}.
\bibitem[IUT4]{IUT4} S.~Mochizuki,
\emph{Inter-universal Teichm\"uller theory IV}, consulted author version,
April 2020. Published in Publ. Res. Inst. Math. Sci. \textbf{57} (2021), 627--723,
\href{https://doi.org/10.4171/PRIMS/57-1-4}{doi:10.4171/PRIMS/57-1-4}.
\bibitem[Cmt]{Cmt} S.~Mochizuki,
\emph{Comments on the manuscript (2018-08 version) by Scholze--Stix concerning
inter-universal Teichm\"uller theory (IUTch)}, September 2018, 5~pp.
\url{https://www.kurims.kyoto-u.ac.jp/~motizuki/Cmt2018-08.pdf}.
\bibitem[Rpt]{Rpt} S.~Mochizuki,
\emph{Report on discussions, held during the period March 15--20, 2018,
concerning inter-universal Teichm\"uller theory (IUTch)},
consulted author version, February 2019, 45~pp.
\url{https://www.kurims.kyoto-u.ac.jp/~motizuki/Rpt2018.pdf}.
\bibitem[JoshiIV]{JoshiIV} K.~Joshi,
\emph{Construction of Arithmetic Teichm\"uller Spaces IV: Proof of the
$abc$-conjecture}, arXiv:2403.10430v2 [math.AG], 24 February 2025;
preliminary version for comments.
\bibitem[First]{First} T.~Urahigashi,
\emph{An explicit initial $\Theta$-datum with mixed reduction at a split prime},
first report, consulted version 0.1.5 (2026),
\href{https://doi.org/10.5281/zenodo.22537342}{doi:10.5281/zenodo.22537342} (all versions).
\bibitem[Second]{Second} T.~Urahigashi,
\emph{An initial $\Theta$-datum over an imaginary quadratic field, and numerical
values of the local quantities of inter-universal Teichm\"uller theory IV},
second report, consulted version 1.0.3 (2026),
\href{https://doi.org/10.5281/zenodo.22601333}{doi:10.5281/zenodo.22601333} (all versions).
\bibitem[Fourth]{Fourth} T.~Urahigashi,
\emph{Split bad primes and the tuple structure of inter-universal
Teichm\"uller theory}, fourth report, consulted version 1.0.2 (2026),
\href{https://doi.org/10.5281/zenodo.22648944}{doi:10.5281/zenodo.22648944} (all versions).
\end{thebibliography}
\end{document}
